\documentclass{article}
\usepackage[left=2cm, right=2cm, top=2cm, bottom=2cm]{geometry}
\usepackage{graphicx}
\usepackage{amsmath}
\usepackage{amssymb}
\usepackage{amsthm}
\usepackage{fancyhdr}
\usepackage{verbatim}
\usepackage{listings}
\usepackage{xcolor}
\usepackage{pdfpages}
\usepackage{cancel}
\usepackage{physics}
\usepackage{esint}
\usepackage{braket}

\lstset{
    language=C++,
    basicstyle=\ttfamily\footnotesize,
    keywordstyle=\color{blue},
    commentstyle=\color{green},
    stringstyle=\color{red},
    numbers=left,
    numberstyle=\tiny\color{gray},
    frame=single,
    breaklines=true,
    captionpos=b,
}


\title{PHYS 487: Lecture \# 22}
\author{Cliff Sun}

\newtheorem{theorem}{Theorem}[section]
\newtheorem{lemma}[theorem]{Lemma}
\newtheorem{definition}[theorem]{Definition}
\newtheorem{conjecture}[theorem]{Conjecture}
\newtheorem{proposition}[theorem]{Proposition}
\newtheorem{corollary}[theorem]{Corollary}
\newtheorem{one minute paper}[theorem]{One Minute Paper}

\pagestyle{fancy}
\lhead{\textbf{\thepage}\ \ \nouppercase{\rightmark}}
\chead{PHYS 487: Lecture \# 22}
\rhead{Cliff Sun}

\begin{document}

\maketitle

\section*{Lecture Span}
\begin{enumerate}
    \item Decoherence
\end{enumerate}

Let the system be called $\ket{S}$ with some spins. Let $\ket{E}$ be the environment be a collection of $N = \mathcal{O}(10^{23})$ other spins. 

Then let the initial state be $$\ket{S} = \frac{\ket{0} + \ket{1}}{\sqrt{2}}$$ Then the total system is 

$$\ket{\psi_{SE}} = \ket{S}_i \otimes \ket{E}_i$$

Let the environment be in states $\ket{e}_0$ and $\ket{e}_1$ (both are combinations of spins in $E$). In principle, the resulting final state $S$ is 

\begin{equation}
    \Tr_{E}[\rho_{SE}] = \rho_s = \sum_i \braket{e_i | \rho_{SE} | e_i} = \sum_i \braket{e_i | \Psi_{SE}}\braket{\Psi_{SE} | e_i}
\end{equation}

This is called a \underline{partial trace}. We would perform this calculation as follows: 
\begin{align}
\frac{1}{2}
\Big(
\bra{e_0}
\left( \ket{0}\ket{e_0} + \ket{1}\ket{e_1} \right)
\left( \bra{0}\bra{e_0} + \bra{1}\bra{e_1} \right)
\ket{e_0}
\;+\;
\bra{e_1}
\left( \ket{0}\ket{e_0} + \ket{1}\ket{e_1} \right)
\left( \bra{0}\bra{e_0} + \bra{1}\bra{e_1} \right)
\ket{e_1}
\Big)
\end{align}

This turns out to be 

\begin{equation}
    \frac{1}{2}\left[ \ket{0}\bra{0} + \ket{1}\bra{1} + \braket{e_1 | e_0}\ket{0}\bra{1} + \braket{e_0|e_1}\ket{1}\bra{0} \right]
\end{equation}

But what makes $\braket{e_0|e_1} = 0$? If there is at least one spin that is not matching between $\ket{e_0}$ and $\ket{e_1}$. Thus, if $\braket{e_1|e_0} \approx 0$, then 

\begin{equation}
    \rho_s = \frac{\ket{0}\bra{0} + \ket{1}\bra{1}}{2}
\end{equation}

This is the basis of decoherence. That is, the environment can change the density matrix of $\rho_s$ which then changes the initial quanutm state. Then, the 
quantum state is no longer in a superposition of $\ket{0}$ and $\ket{1}$, but rather only in either $\ket{0}$ or $\ket{1}$.  

As a starting point, 

\begin{equation}
    \dot{\rho}_{SE} = -\frac{i}{\hbar}\left[ H_{SE}, \rho_{SE} \right]
\end{equation}

Then 

\begin{equation}
    \rho_{SE}(t + \Delta t) = u_{SE}(\Delta t)\left[ \rho_{s}(t) \otimes \rho_E(t) \right]u^{\dagger}_{SE}(\Delta t)
\end{equation}

Let there be a "jump operator" 

\begin{equation}
    L_n = \braket{e_n | u | e_0}
\end{equation}

Which dictates the probability of the environment jumping from $e_0$ to $e_n$. Then 

\begin{equation}
    \rho_s(t + \Delta t) = \sum_n L_n \rho_s(t)L_{n}^{\dagger}
\end{equation}

We can then expand 

\begin{equation}
    L_0 \approx \mathbb{I} + \left( -iH_s - \frac{1}{2} \sum_n A_n^{\dagger}A_n\right)dt
\end{equation}

Where 

\begin{equation}
    L_n \approx \sqrt{\gamma_n}dt A_n
\end{equation}

Here, $\gamma_n$ is the rate of evolution of $E$. Then, you can rewrite everything into a differential equation for $\rho_s$. That is 

\begin{equation}
    \dot{\rho_s} = -i [H_s, \rho_s] + \sum_n \gamma_n (A_n \rho_s A_n^{\dagger}) - \frac{1}{2}\{ A_n^{\dagger}A_n, \rho_s \}
\end{equation}

\section*{Measurement -- Von Neumann Model}

Suppose there is a pointer $\ket{P_i}$. Then we want 

\begin{equation}
    \alpha\ket{0} + \beta\ket{1} \rightarrow \alpha\ket{0}\ket{p_0} + \beta\ket{1}\ket{p_1}
\end{equation}

This is also called a "cat state". Since measuring if a particle decayed means that the cat is dead, etc. 

After applying some time evolution operator of 

$$u = \exp(-i H_{int}t/\hbar)$$

Then 

\begin{equation}
    H_{int} = \hbar g_{int}\sigma_z \otimes p_i
\end{equation}



\end{document}